%% 第26回編集委員会（物理第4巻）：事務局
%% 物理シリーズ　第4巻　電磁気学　巻末（こたえ・記号カード・物性値・定数・単位・なぜ？の道すじ・チェック表・つぎの巻へ・さくいん・奥付）
%% 第3巻の ph3-back.tex と同じ部品・同じ形。見出し（\bkbacktitle）の名前は、もくじ（ph4-toc.tex）の巻末の項目の名前と同じにする
%%   （tools/check_pages.py p4 が照らし合わせる。表の名前は phys/series_phys.py の第4巻の tables「物性値・定数・単位」）。
%% 第1〜3巻の確認で決めたこと（U-1〜U-13、Q6、X-1〜X-8）を、はじめから当てはめた：
%%   ・記号カード：数学の記号の行は、この巻の章の \bkkigocard（自動でうしろに集まる）と同じ記号・同じ意味を置かない（U-6・X-2）。
%%   ・紙面に委員会の回を書かない（U-7）。回はこのコメントにだけ書く。
%%   ・さくいんの段は \raggedcolumns。奥付の「先に読む本」は表4（bk-phys.sty）と同じ文（Q6-22）。
%%   ・「分野で記号がぶつかるもの」は、X-8 のやり方（同じ記号を2行にしない、章で確かめた意味だけを書く）で、11月5日の15章の原稿の
%%     頭のコメント（記号の欄）と本文で確かめて作った。
%% 物性値の出典は、各章のファイルの頭のコメント（出典の欄）と docs/作業メモ/物理第4巻_組どうしの確かめ.md の「巻末の表に入れる値」
%%   「まとめ役が決めたこと」Z-1〜Z-9 から集めた。計算は checks/物理第4巻/巻末/calc_ph4_back.py（出力 calc_ph4_back_out.txt）。
%%   ・定数（ε0・μ0・e・c・mₑ・mₚ・F・真空のインピーダンス 376.730 313 412 Ω）：NIST の CODATA の値のページを、11月5日に事務局が curl で読み直した。
%%   ・銅の抵抗率は Z-4（標準軟銅 1.72×10⁻⁸、導電率 5.80×10⁷ S/m）。Pt100・K 熱電対は規格の式で計算し直して、電-4 の値と同じになった。
%%   ・地磁気は Z-6（水平 3.0×10⁻⁵ T・鉛直 3.5×10⁻⁵ T・伏角 約49°。〔要確認〕のまま）。同軸は Z-5（0.45 mm・1.6 mm）。
%%   ・ほかの値は、各章の組が確かめた出典（章の頭のコメント）のまま。Wikipedia は出典にしていない。
%%   〔要確認〕は、一次資料でも、一次資料を引く2つ以上の資料でも確かめられなかった値（章の頭のコメントで〔要確認〕のもの）。
%%     比誘電率（空気・ポリエチレン・PET・ガラス・チタン酸バリウム）、黒鉛・ガラス・ポリエチレンの抵抗率、鉄・電磁鋼板の比透磁率、
%%     地磁気、銅の自由電子の数。ポリエチレン 2.3 は、電-13 が RG-58 のデータシートの速度係数 0.66（εr＝1/0.66²＝2.30）と電子部品の
%%     資料の表で照らし合わせているが、作業メモが〔要確認〕としているので、そのままにした（委員会で外すかを決める）。
%% 表の書式（見出しの行の色・横の線・行の間）は、第3巻と同じく、各表の colspec のあとにじかに書いた（マクロにまとめると tabularray が展開しない）。
\makeatletter
\def\phbkhead#1{\par\addvspace{4mm}\Needspace{6\baselineskip}\noindent{\sffamily\bfseries\large\color{bkDeep}#1}\par\nobreak\vspace{2mm}}
\newcommand\phyk{〔要確認〕}
%% ---------------- こたえ ----------------
\bkAnswerPages
\par\vspace{4mm}
%% 箱はページをまたがせない（またぐと、次のページに箱の切れはしだけが残る）
\begin{tcolorbox}[bk base, unbreakable, colframe=bkTheme!60, colback=bkTheme!5, boxrule=0.8pt, arc=2mm, left=4mm, right=4mm, top=2mm, bottom=2mm]
{\sffamily\bfseries\color{bkDeep}答え合わせのしかた}\par\vspace{1mm}
\begin{enumerate}[label={\bkcircled{\arabic*}}, leftmargin=2.4\zw, itemsep=0pt, topsep=0pt]
\item まず自分の答えを全部書いてから、ここを見よう。
\item 単位と桁を確かめよう。$10^{-12}$ や $10^{9}$ の指数、$f$〔Hz〕と $\omega=2\pi f$〔rad/s〕、交流の実効値と山の値（$\sqrt2$ 倍）、W と var と VA を取り違えていないかも。
\item 丸め方で最後の桁が1つずれても、間違いではありません。$k=8.99\times10^9$ と $1/(4\pi\varepsilon_0)$ のどちらで計算しても、2けたは同じになります。
\item 違ったら、同じ形の例題にもどって、どの段で違ったかを探そう。$\phi$ が電位か位相か、$e$ が電気素量か電源の電圧かも確かめよう。
\end{enumerate}
\end{tcolorbox}

%% ---------------- 記号カード ----------------
%% 数学シリーズの記号で、この巻でよく使うもの（この巻の章の \bkkigocard は自動でうしろに集まる）。
%% 章のカードにある記号（11月5日の原稿で数えた。電-1：Q・q、E、k、ε0、e、λ・σ、p　電-2：φ、V、Φ_E、ρ　電-3：C、ε、D、P
%%   電-4：I、R、ρ、σ、P、𝓔　電-5：G、R_th　電-6：B、μ0　電-7：H、M、μ・μr、𝓡_m　電-8：L、M、δ　電-9：e、v、τ、ζ、s
%%   電-10：Z、X　電-11：Q、S、V_p、n_s　電-12：T、K、s、J　電-13：Z0、L'、γ、r_V　電-14：S、I_d　電-15：η、λ_c）は、同じ意味ではここに置かない（U-6）。
%%   ラプラス変換の s は電-9 のカードにあるので、9-5 の行は 𝓛 だけにした。φ は電位（電-2 のカード）と意味が違う「位相」の行を置いた（X-2）。
\bkKigoPages{%
\bkkigorow{$\Delta$}{デルタ}{変化量・きざみ。小さな長さ $\Delta l$、時間のきざみ $\Delta t$}{2-3}%
\bkkigorow{$\omega$}{オメガ}{角速度・角周波数〔rad/s〕。$\omega=2\pi f$}{4-8}%
\bkkigorow{$\vec a\cdot\vec b$}{エー・ドット・ビー}{内積。$\lvert\vec a\rvert\lvert\vec b\rvert\cos\theta$。答えは数。電場の流束の $\vec E\cdot\vec n$、仕事}{4-12}%
\bkkigorow{$\vec a\times\vec b$}{エー・クロス・ビー}{外積。ローレンツ力 $q\vec v\times\vec B$、ポインティングベクトル $\vec E\times\vec H$}{4-14}%
\bkkigorow{$\odot$、\ $\otimes$}{手前へ・おくへ}{紙に垂直な向き（電流・磁場の向き）。$\odot$ は手前へ出る、$\otimes$ はおくへ入る}{4-14}%
\bkkigorow{$\dot x$、\ $\ddot x$}{エックス・ドット、エックス・ツードット}{時間で1回微分・2回微分。電流 $i=\dot q$}{5-5}%
\bkkigorow{$\approx$}{ニアリーイコール}{ほぼ等しい（近似）}{5-8}%
\bkkigorow{$\displaystyle\int$、\ $\displaystyle\sum$}{インテグラル、シグマ}{細かく分けてたす・全部たす。広がった電荷の電場、平均の電力}{5-10}%
\bkkigorow{$e^x$、\ $\ln x$}{イー、エル・エヌ・エックス}{ネイピア数（2.718…）と自然対数。$e^{-t/\tau}$。電気素量・電源の電圧の $e$ とは別}{5-14}%
\bkkigorow{$\phi$（位相）}{ファイ}{位相のずれ・力率の角（電-10・11）。$\cos\phi$。電位 $\phi$（電-2）とは別の意味}{5-15}%
\bkkigorow{$\log_{10}x$、\ dB}{ログ、デシベル}{常用対数と、比を「10×常用対数」で表す単位。電力の比は $10\log_{10}$、電圧の比は $20\log_{10}$}{3-12}%
\bkkigorow{$j$}{ジェイ}{虚数単位（2乗すると $-1$）。電流の $i$ とぶつけないため $j$ と書く}{6-8}%
\bkkigorow{$\lvert z\rvert$、\ $\arg z$}{絶対値、アーギュメント}{複素数の大きさと偏角。フェーザの大きさ（実効値）と位相}{6-9}%
\bkkigorow{$e^{j\theta}$}{イー・ジェイ・シータ}{単位円の上の角 $\theta$ の点（オイラーの公式）。回る矢印（フェーザ）}{6-10}%
\bkkigorow{$\cosh x$、\ $\sinh x$}{ハイパボリック・コサイン、ハイパボリック・サイン}{双曲線関数。損失のある長い線の電圧・電流（電-13）}{6-11}%
\bkkigorow{$\det$}{デターミナント}{行列式。四端子定数の $AD-BC=1$（電-13）}{7-4}%
\bkkigorow{$\mathbf A^{-1}$}{エー・インバース}{逆行列。網目電流法・節点電位法の連立方程式（電-5）}{7-5}%
\bkkigorow{$\tfrac{\partial f}{\partial x}$}{ラウンドディー}{偏微分。電信方程式・波動方程式}{8-2}%
\bkkigorow{$\nabla$、\ $\mathrm{grad}$}{ナブラ、グラッド}{勾配。$\vec E=-\mathrm{grad}\,\phi$（電-2）}{8-4}%
\bkkigorow{$\displaystyle\oint_C$}{周回積分}{閉じた道を1周する線積分。アンペールの法則、ファラデーの法則}{8-7}%
\bkkigorow{$\displaystyle\oiint_S$}{閉じた面の面積分}{閉じた袋の面でたす。ガウスの法則。面のタイルは $\vec n\,\dd S$（電-14・15 は $\dd A$）}{8-8}%
\bkkigorow{$\mathrm{div}$}{ダイバージェンス}{発散。$\mathrm{div}\,\vec E=\rho/\varepsilon_0$、$\mathrm{div}\,\vec B=0$}{8-9}%
\bkkigorow{$\mathrm{rot}$}{ローテーション}{回転（うず）。$\mathrm{rot}\,\vec B=\mu_0\vec j$、$\mathrm{rot}\,\vec E=-\partial\vec B/\partial t$}{8-10}%
\bkkigorow{$\nabla^2$}{ラプラシアン}{ポアソン方程式・ラプラス方程式（電-2）、波動方程式（電-15）}{8-12}%
\bkkigorow{$\mathcal L$}{ラプラス変換}{微分が「$s$ をかける」になる（電-9。$s$ のカードは電-9）}{9-5}%
}

\Needspace{18\baselineskip}\phbkhead{分野で記号がぶつかるもの}
{\small 電磁気学は、場・回路・機械・電波と分野が広く、同じ文字が違う量を表すことがあります。この本では、同じ章の中では1つの文字を1つの意味にだけ使い、ぶつかるときは文字や添字・書体を変えました。どの量かは、章（部）と単位で見分けよう。右から2つめの列は、第1巻『力学』・第2巻『熱学』・第3巻『波動』とほかの巻での意味です。\par}\vspace{2mm}
{\footnotesize
\begin{longtblr}{colspec={Q[c,m,wd=15mm] X[l,m] Q[l,m,wd=34mm] Q[l,m,wd=28mm]}, width=\linewidth, rowhead=1, row{1}={bg=bkVol!14, font=\sffamily\bfseries\footnotesize}, hlines={0.4pt,bkLine}, rowsep=1.4pt, colsep=3.5pt}
記号 & この巻での意味 & 第1〜3巻・ほかの巻での意味 & 見分け方 \\
$c$ & 真空の光速〔m/s〕（電-0・14・15）、線の上を進む波の速さ（電-13。真空の光速は $c_0$）、相似の表のダンパーの係数（前付・電-9 ⑨） & ダンパーの係数（力-2・波-1〜4）、比熱（熱-1）、波の速さ（波-5〜11） & 部で見分ける。単位：m/s \\
$S$ & 面積〔m$^2$〕（電-1〜13）、皮相電力〔VA〕（電-11。面積を使わない）、ポインティングベクトル $\vec S$〔W/m$^2$〕（電-14・15） & エントロピー（熱-6）、断面積（波-1〜9）、安全率（力-8） & 単位：m$^2$、VA、W/m$^2$ \\
$A$ & 面積〔m$^2$〕（第4部：電-14・15、電-16）、四端子定数の $A$（電-13 の行列 $\begin{pmatrix}A&B\\C&D\end{pmatrix}$。単位なし） & 振幅（波）、断面積（力-8） & 第4部と電-16 は面積。行列の成分かどうか \\
$R$ & 電気抵抗〔Ω〕（電-4〜13）、円板・球・極板の半径（電-1・2・14）、反射率（電-15。放射抵抗は $R_{\mathrm r}$）。テブナンの等価抵抗 $R_{\mathrm{th}}$（電-5）。磁気抵抗は書体を変えて $\mathcal R_{\mathrm m}$（電-7） & 気体定数（熱-3）、熱抵抗 $R_{\mathrm{th}}$（熱-11・前付）、曲率半径（波-12） & $R_{\mathrm{th}}$ は単位で：Ω と K/W \\
$L$ & 自己インダクタンス〔H〕（電-8〜13）、1\,m あたり $L'$〔H/m〕（電-13）、網の厚さ（電-15 だけ）。長さは $l$ & ラグランジアン（波-3）、潜熱（熱-1）、長さ（力-8）、角運動量 $\vec L$（力-6） & 単位：H、H/m \\
$C$ & 静電容量〔F〕（電-3・5・9〜13）、1\,m あたり $C'$〔F/m〕（電-13）、四端子定数の $C$（電-13）。単位のクーロン C は立てた字 & 熱容量 $C_{\mathrm{th}}$（熱-1）、通信路容量（波-11） & 単位：F、F/m \\
$Q$、$q$ & 電荷〔C〕（全章。場をつくる・たまった $Q$、試しの・小さな $q$）、Q 値（電-10）、無効電力〔var〕（電-11。電荷を使わない）、電-4 の熱は $Q_{\text{熱}}$ & 熱量（熱-1）、体積流量（力-13）、Q 値（波-2）、一般化座標 $q$（波-3） & 単位：C、var。Q 値は単位なし \\
$P$、$p$ & 電力〔W〕（電-4・11）、分極 $\vec P$〔C/m$^2$〕（電-3）、電磁波を出す仕事率（電-15）。双極子モーメント $\vec p$（電-1〜3）、瞬時電力 $p$（電-11）、極の数 $p$（電-12）、放射圧 $p$（電-15） & 仕事率（力-4）、圧力 $p$（熱）、運動量 $\vec p$（力-3） & 矢印の有無と単位：W、C/m$^2$、C·m、Pa \\
$T$ & トルク〔N·m〕（電-6 ⑪・電-12。トルクの章では周期に使わない）、周期〔s〕（電-10） & 周期（力-1）、ねじりモーメント（力-11）、絶対温度（熱-1）。トルクは $\vec N$（力-6）・$\tau_{\mathrm m}$（波-4） & 単位：N·m、s \\
$E$、$\mathcal E$ & 電場 $\vec E$〔V/m〕（第1・2・4部）、電源の電圧（直流・実効値）$E$〔V〕（電-9〜13）、逆起電力 $E$（電-12）。場の章の起電力は書体を変えて $\mathcal E$（電-4・5・8） & ヤング率（力-8）、エネルギー（力-4）、照度（波-12） & 電場は V/m、電圧は V \\
$e$ & 電気素量〔C〕（電-1・2・4・6）、電源の電圧の瞬時値〔V〕（電-9〜13）。肩に式がのった $e^{x}$ はネイピア数 & 反発係数（力-3）、ネイピア数（5-14） & 単位：C、V \\
$V$、$v$ & 電圧（電位の差）$V$〔V〕、部品の電圧の瞬時値 $v$（電-9〜13）、速さ $v$（電-1・4・6・8・14。ドリフトの速さ $v_{\mathrm d}$）。単位のボルト V は立てた字 & 体積 $V$（力・熱）、速度 $v$（力-1）、比体積 $v$（熱-8） & 単位：V、m/s \\
$\phi$、$\Phi$ & 電位 $\phi$〔V〕（電-1〜8・14）、位相のずれ $\phi$（電-10〜13。電位を使わない章）。磁束 $\Phi$〔Wb〕（電-7・8）、電場の流束 $\Phi_E$（電-2・14）・磁束 $\Phi_B$（電-14） & 初期位相（波-1）、相対湿度 $\varphi$（熱-9） & 単位：V、rad、Wb \\
$B$ & 磁場（磁束密度）$\vec B$〔T〕（電-6〜8・14・15）、四端子定数の $B$〔Ω〕（電-13）。サセプタンスは記号にしない & 高校の教科書の「磁束密度」 & 単位：T \\
$H$ & 磁場の強さ $\vec H$〔A/m〕（電-7・14・15。高校の教科書の「磁場」）、発電機の慣性定数〔s〕（電-16） & エンタルピー（熱-4）、全ヘッド（力-13）、ゆれの倍率・情報のエントロピー（波-10・11） & 単位：A/m \\
$M$ & 磁化 $\vec M$〔A/m〕（電-7）、相互インダクタンス $M$〔H〕（電-8・12） & 平面のモーメント（力-5）、モル質量（熱-3）、質量行列 $\mathbf M$（波-3） & 章と単位で：A/m、H \\
$\rho$ & 電荷密度〔C/m$^3$〕（電-2・14。電-4 は $\rho_{\mathrm q}$）、抵抗率〔Ω·m〕（電-4・8）、密度〔kg/m$^3$〕（電-7 例題10） & 密度（力-12）、反射率（熱-13） & 単位：C/m$^3$、Ω·m \\
$\sigma$ & 面電荷密度〔C/m$^2$〕（電-1〜3）、導電率〔S/m〕（電-4・8・14・15） & 応力（力-8）、ステファン・ボルツマン定数（熱-13） & 単位：C/m$^2$、S/m \\
$\varepsilon$ & 誘電率〔F/m〕（電-3。真空は $\varepsilon_0$、比誘電率 $\varepsilon_{\mathrm r}$）、ひずみ（電-4 ⑧のひずみゲージ） & ひずみ（力-8）、放射率・圧縮比（熱-13・熱-7） & 誘電率は F/m、ひずみは単位なし \\
$\mu$ & 透磁率〔N/A$^2$〕（電-7。真空は $\mu_0$、比透磁率 $\mu_{\mathrm r}$）。単位の μ（マイクロ）は立てた字 & 摩擦係数・粘度（力-2・力-14）、線密度（波-5）、化学ポテンシャル（熱-10） & 単位：N/A$^2$＝H/m \\
$\eta$ & 波のインピーダンス〔Ω〕（電-15。真空 $\eta_0\approx\qty{377}{\ohm}$）、効率（電-12・16）、浅い水の波の水面の高さ（電-13 の表） & 効率（力-4・熱-5）、水面の高さ（波-5） & 単位：Ω、効率は単位なし \\
$Z$、$X$ & インピーダンス $Z$〔Ω〕（電-10〜13）、特性インピーダンス $Z_0$（電-13）、リアクタンス $X$〔Ω〕（電-10〜13） & 波のインピーダンス $Z$（波-7）、圧縮係数（熱-3）、複素振幅 $X$（6-13・波-2） & 単位：Ω \\
$k$、$K$ & クーロンの定数 $k$（電-1〜3）、波数 $k$（電-15。クーロンの定数は $1/(4\pi\varepsilon_0)$）、ばね定数 $k$（相似の表）。ゲージ率 $K$（電-4）、トルク定数 $K$（電-12） & ばね定数（力-2）、熱伝導率（熱-11）、波数（波-5）、ゲイン $K$（波-4） & 単位：N·m$^2$/C$^2$、rad/m、N·m/A \\
$s$ & ラプラス変換の変数（電-9）、すべり（電-12。単位なし） & ラプラス変換（9-5・波-4）、比エントロピー（熱-6） & 章で見分ける \\
$I$、$J$ & 電流 $I$〔A〕（瞬時値は $i$）、電磁波の強さ $I$〔W/m$^2$〕（電-15。電流は記号にしない）。電-12 の慣性モーメントは $J$ & 慣性モーメント $I$（力-6）、波の強さ（波-7） & 単位：A、W/m$^2$ \\
$n$ & 1\,m あたりの巻数（電-6〜8）、自由電子の数密度（電-4）、回転数〔min$^{-1}$〕（電-12）、屈折率（電-15） & 物質量（熱-3）、屈折率（波-12） & 単位：1/m、1/m$^3$ \\
$\delta$ & 表皮深さ〔m〕（電-8・15） & 対数減衰率（波-1）、デルタ関数（9-6） & 単位：m \\
$\lambda$ & 線電荷密度〔C/m〕（電-1・3）、波長〔m〕（電-13・15） & 波長（波-5）、管摩擦係数（力-14） & 単位：C/m、m \\
$\tau$ & 時定数〔s〕（電-9）、電子が衝突する間の平均の時間（電-4） & せん断応力（力-8）、透過係数（波-7）、モーターのトルク $\tau_{\mathrm m}$（波-4） & 単位：s \\
$\alpha$、$\gamma$ & 抵抗の温度係数 $\alpha$〔1/K〕（電-4）、減衰定数 $\alpha$・伝搬定数 $\gamma$（電-13） & 線膨張係数（力-8）、比熱比 $\gamma$（熱-3） & 単位：1/K、1/m \\
$G$ & コンダクタンス〔S〕（電-5。節点の行列は $\mathbf G$）、1\,m あたり $G'$（電-13） & 横弾性係数・万有引力定数（力-8・力-2）、ギブズの自由エネルギー（熱-10）、伝達関数 $G(s)$（9-8） & 単位：S \\
$U$、$u$ & 静電エネルギー・位置エネルギー $U$〔J〕（電-2・3・8）、エネルギー密度 $u$〔J/m$^3$〕（電-3・8・14） & 内部エネルギー（熱-2）、位置エネルギー（力-4） & 単位：J、J/m$^3$ \\
\end{longtblr}\par}

\phbkhead{この巻の約束（記号と向き）}
{\small
\begin{itemize}[leftmargin=1.2\zw, itemsep=2pt, topsep=1pt, label={\color{bkVol}●}]
\item \textbf{磁場}：$\vec B$ を「磁場（磁束密度）」とよびます（数学 8-10・8-14 と同じ）。$\vec H$ は電-7 から「磁場の強さ」です。高校の教科書の「磁場 $H$」と「磁束密度 $B$」は、真空で $B=\mu_0H$ の関係です（電-6 ③）。専門書の「電界・磁界」は、電場・磁場と同じものです。
\item \textbf{電圧の書き方}：電位は $\phi$、電圧（2点の電位の差）は $V$、場の章の起電力は $\mathcal E$。回路の章（電-9〜13）では、電源の電圧を瞬時値 $e$・直流と実効値 $E$、部品の電圧を瞬時値 $v$・実効値 $V$、電流を瞬時値 $i$・実効値 $I$ と書きます（前付の相似の表、数学 6-15 と同じ）。
\item \textbf{交流}：ことわりがなければ、電圧・電流は実効値です（100\,V の山は約141\,V）。周波数 $f$〔Hz〕と角周波数 $\omega$〔rad/s〕を書き分けます。第3巻の「振動数」と同じ量です。
\item \textbf{面積}：電-1〜13 は $S$、第4部（電-14・15）は $A$ です（ポインティングベクトル $\vec S$ とぶつけない）。
\item \textbf{向き}：電流は正の電荷が動く向き（電子が動く向きの逆）です。電流と磁場の向きは右ねじの法則で決め（電-6 ②）、紙に垂直な向きは $\odot$（手前へ）・$\otimes$（おくへ）で書きます。誘導起電力の向きはレンツの法則で決めます（電-8 ③）。
\item \textbf{反射係数}：電-13 の電圧の反射係数 $r_V$ は、波-7 の力・音圧で数えた反射係数 $r_{\text{押}}$ と同じ式、電流の反射係数 $r_I=-r_V$ は波-7 の $r$ と同じ式です。電-15 の電場の反射係数 $r_E$ も $r_{\text{押}}$ と同じ式です。
\end{itemize}\par}

%% ---------------- 物性値・定数・単位 ----------------
\bkbacktitle{物性値・定数・単位}
{\small この巻で使った値です。本文の例題の値は、この表と同じ丸めです。物性値は温度や材料の種類で変わるので、設計では便覧・規格の値を使います。出典の欄の「電-4」などは、その値を使った章です。\phyk は、規格・公式資料の本文でも、それを引く2つ以上の資料でも、まだ確かめられていない値です。\par}

\phbkhead{1　抵抗率と導電率（\qty{20}{\celsius}）}
{\footnotesize
\begin{longtblr}{colspec={Q[l,m,wd=34mm] Q[c,m,wd=27mm] Q[c,m,wd=22mm] X[l,m]}, width=\linewidth, rowhead=1, row{1}={bg=bkVol!14, font=\sffamily\bfseries\footnotesize}, hlines={0.4pt,bkLine}, rowsep=1.4pt, colsep=3.5pt}
物 & {抵抗率 $\rho$\\〔Ω·m〕} & {導電率 $\sigma$\\〔S/m〕} & 出典・注意 \\
銀 & $1.59\times10^{-8}$ & $6.3\times10^{7}$ & ア（電-4） \\
銅（電線の標準の銅） & $1.72\times10^{-8}$ & $5.80\times10^{7}$ & イ（電-4・8・13・15） \\
アルミニウム & $2.65\times10^{-8}$ & $3.8\times10^{7}$ & ア（電-4・8） \\
タングステン & $5.6\times10^{-8}$ & $1.8\times10^{7}$ & ア（電-4・9） \\
鉄 & $9.7\times10^{-8}$ & $1.0\times10^{7}$ & ア（電-4） \\
ニクロム（電熱線） & $1.1\times10^{-6}$ & $9\times10^{5}$ & ア（電-4） \\
黒鉛（鉛筆の芯） & 約 $10^{-5}$ & 約 $10^{5}$ & 目安\phyk（電-4） \\
海水（塩分35、\qty{15}{\celsius}） & 0.233 & 4.29 & ウ（電-14） \\
ガラス & $10^{10}$〜$10^{14}$ & ― & 目安\phyk（電-4） \\
ポリエチレン & $10^{16}$ 以上 & ― & 目安\phyk（電-4） \\
\end{longtblr}\par}
{\footnotesize\color{bkGray}出典　ア：物理の教科書の表（Giancoli の値）を引く大学の資料2つを照らし合わせた値。イ：焼きなました銅の国際標準（IEC 60028。$\tfrac1{58}$\,Ω·mm$^2$/m＝$1.724\times10^{-8}$\,Ω·m、導電率 $5.80\times10^7$\,S/m を 100\,\% とする）。米国標準局（NBS）の銅線表で照らし合わせた。ウ：標準海水の導電率 42.914\,mS/cm（TEOS-10、Culkin・Smith 1980）。長さ 1\,m・断面積 1\,mm$^2$ の線の抵抗は、$\rho$〔Ω·m〕の $10^6$ 倍（銅 0.017\,Ω、ニクロム 1.1\,Ω。電-4 ③）。\par}

\phbkhead{2　抵抗の温度係数と、温度をはかる素子}
{\footnotesize
\begin{longtblr}{colspec={Q[l,m,wd=40mm] Q[l,m,wd=46mm] X[l,m]}, width=\linewidth, rowhead=1, row{1}={bg=bkVol!14, font=\sffamily\bfseries\footnotesize}, hlines={0.4pt,bkLine}, rowsep=1.4pt, colsep=3.5pt}
名前 & 値 & 出典・注意 \\
銅の抵抗の温度係数 & 0.00393\,/K（\qty{20}{\celsius}） & IEC 60028・米国標準局の銅線表（電-4） \\
タングステンの温度係数 & 約 0.0045\,/K（\qty{20}{\celsius}付近） & 高温では一定でない。電球の熱いときと冷たいときの抵抗の比は 10〜15 ほど（電-4・9） \\
白金の測温抵抗体 Pt100 & \qty{0}{\celsius}で\qty{100.00}{\ohm}、\qty{100}{\celsius}で\qty{138.51}{\ohm}（$\alpha=0.00385$\,/K） & JIS C 1604:2013・IEC 60751:2008 の式で計算し直した（電-4） \\
K 熱電対（基準接点 \qty{0}{\celsius}） & \qty{20}{\celsius}で\qty{0.798}{mV}、\qty{300}{\celsius}で\qty{12.209}{mV}（平均 約\qty{41}{\micro V/K}） & NIST Monograph 175 の ITS-90 の式で計算し直した（電-4。熱-1） \\
燃料電池（水素と酸素） & 起電力 \qty{1.23}{V}（液体の水ができるとき） & $\Delta G=-\qty{237.13}{kJ/mol}$（熱-10）から（電-4） \\
\end{longtblr}\par}

\phbkhead{3　比誘電率}
{\footnotesize
\begin{longtblr}{colspec={Q[l,m,wd=44mm] Q[c,m,wd=26mm] X[l,m]}, width=\linewidth, rowhead=1, row{1}={bg=bkVol!14, font=\sffamily\bfseries\footnotesize}, hlines={0.4pt,bkLine}, rowsep=1.4pt, colsep=3.5pt}
物 & 比誘電率 $\varepsilon_{\mathrm r}$ & 出典・注意 \\
真空 & 1 & 定義 \\
空気（1気圧） & 約1.0006 & 目安\phyk（電-3） \\
ポリエチレン & 約2.3 & 同軸ケーブル（RG-58）のデータシートの速度係数 0.66 から $1/0.66^2=2.30$。\phyk（電-3・13） \\
発泡ポリエチレン & 約1.5（1.5〜1.6） & テレビの同軸ケーブル（5C-FB）。ケーブルの会社のカタログ・仕様書（電-13） \\
PET（クリアファイルなど） & 約3 & 資料で 3.0〜3.3 とばらつく\phyk（電-3） \\
ガラス & 約4〜10 & 種類で大きく違う\phyk（電-3） \\
水（\qty{20}{\celsius}・\qty{25}{\celsius}） & 80.2・78.4 & IAPWS の水の誘電率の式（R8-97）で計算（電-1・3・14・15） \\
チタン酸バリウム（セラミックス） & 1000 以上 & 小さなコンデンサーの材料\phyk（電-3） \\
\end{longtblr}\par}
{\footnotesize\color{bkGray}比誘電率は、電場の向きが変わる速さ（周波数）と温度で変わります。水の 80 は、ゆっくり変わる電場での値で、光（約 $10^{15}$\,Hz）では屈折率の2乗の約1.8 になります（電-15 ⑨）。同軸ケーブルの中を進む電気の波の速さは、光速の $1/\sqrt{\varepsilon_{\mathrm r}}$ 倍（ポリエチレンで 0.66 倍、発泡ポリエチレンで約 0.8 倍。電-13 ②）。\par}

\phbkhead{4　磁性と磁石}
{\footnotesize
\begin{longtblr}{colspec={Q[l,m,wd=40mm] Q[l,m,wd=46mm] X[l,m]}, width=\linewidth, rowhead=1, row{1}={bg=bkVol!14, font=\sffamily\bfseries\footnotesize}, hlines={0.4pt,bkLine}, rowsep=1.4pt, colsep=3.5pt}
名前 & 値 & 出典・注意 \\
磁化率 $\chi$（体積あたり、SI） & 水 $-9.0\times10^{-6}$、銅 約 $-1.0\times10^{-5}$、アルミニウム $+2.2\times10^{-5}$ & 反磁性浮上の論文（arXiv:2004.09683 ほか）、磁性の用語集・講義ノート（CRC Handbook を引く）（電-7） \\
鉄・電磁鋼板の比透磁率 & 数百〜数千（磁場の強さで変わる） & 材料と磁場の強さで大きく変わる\phyk。例題の値は仮の値（電-7・8・12） \\
鉄のキュリー温度 & \qty{770}{\celsius}（\qty{1043}{K}） & 大学の材料科学の演示実験の資料、分析機器の会社の用語集（電-7） \\
電磁鋼板の飽和 & 約\qty{2}{T} & 鉄鋼の会社の製品の説明、方向性電磁鋼板の特許（電-7・12） \\
電磁鋼板の鉄損（35A300） & \qty{1.5}{T}・\qty{50}{Hz}で\qty{3.00}{W/kg}以下 & JIS C 2552:2014（電-7） \\
磁石の残留磁束密度 & ネオジム磁石 約\qtyrange{1.2}{1.4}{T}、フェライト磁石 約\qty{0.4}{T} & 部品の会社の技術資料・磁石の会社のカタログ（電-7） \\
身の回りの磁場 & 冷蔵庫の磁石の表面 約\qty{0.01}{T}、MRI \qtyrange{1.5}{3}{T} & 米国国立強磁場研究所（MagLab）の教材（電-6） \\
地磁気（東京の近く） & 水平 約$3.0\times10^{-5}$\,T、鉛直 約$3.5\times10^{-5}$\,T、全体 約$4.6\times10^{-5}$\,T、伏角 約\ang{49}、偏角 約\ang{7}西 & IGRF で計算した値を引く大学の授業資料。数学 8-10 の「日本ではおよそ $4.5\times10^{-5}$\,T」と合う\phyk（電-6・7・8） \\
\end{longtblr}\par}

\phbkhead{5　大気と雷・静電気}
{\footnotesize
\begin{longtblr}{colspec={Q[l,m,wd=40mm] Q[l,m,wd=46mm] X[l,m]}, width=\linewidth, rowhead=1, row{1}={bg=bkVol!14, font=\sffamily\bfseries\footnotesize}, hlines={0.4pt,bkLine}, rowsep=1.4pt, colsep=3.5pt}
名前 & 値 & 出典・注意 \\
晴れた日の地面の近くの電場 & 約\qty{100}{V/m}（下向き） & 地球全体の電気回路の総説（Siingh ほか 2007）、米国科学アカデミーの報告（1986）（電-2） \\
電離層と地面の電位の差 & 約\qty{300}{kV} & 同じ総説（電-2） \\
空気がこわれる電場（絶縁破壊） & 約 $3\times10^6$\,V/m（1気圧・常温・平らな電極） & 高電圧工学の講義資料と解説で照らし合わせた（電-2・3） \\
雷1回で地面へ運ばれる電荷 & 数\,C から数十\,C & 国際大電力システム会議（CIGRE）の雷のパラメータを引く資料、雷の教科書（電-1） \\
人がパチッと感じる体の電圧 & 約\qty{3}{kV} & 帯電の研究の報告、技術資料（人の体の静電容量 約\qty{100}{pF}）（電-1） \\
雷のときの車の中 & 比較的安全（オープンカーは除く） & 気象庁「雷から身を守るには」（電-3） \\
\end{longtblr}\par}

\phbkhead{6　法令で決まっている値と、電波の周波数}
{\footnotesize
\begin{longtblr}{colspec={Q[l,m,wd=40mm] Q[l,m,wd=46mm] X[l,m]}, width=\linewidth, rowhead=1, row{1}={bg=bkVol!14, font=\sffamily\bfseries\footnotesize}, hlines={0.4pt,bkLine}, rowsep=1.4pt, colsep=3.5pt}
名前 & 値 & 出典・注意 \\
家庭の電圧 & 100\,V の所は 101\,V の上下 6\,V 以内、200\,V の所は 202\,V の上下 20\,V 以内 & 電気事業法施行規則 第38条（e-Gov 法令）（電-4・10・11） \\
送電線などの近くの磁場 & 磁束密度 \qty{200}{\micro T}以下（商用周波数） & 電気設備に関する技術基準を定める省令 第27条の2（e-Gov 法令）（電-6・11） \\
電波 & 300万\,MHz（\qty{3}{THz}）以下の周波数の電磁波 & 電波法 第2条（e-Gov 法令）（電-15） \\
電磁誘導加熱式調理器（IH） & \qtyrange{20.05}{100}{kHz} & 電波法施行規則 第46条の7（e-Gov 法令）（電-8・15） \\
電子レンジ & \qty{2450}{MHz}±\qty{50}{MHz}（波長 約\qty{12.2}{cm}） & 同じ（電-15。波-6） \\
IC カード（読み書きの装置） & \qty{13.56}{MHz} & 電波法施行規則 第44条（e-Gov 法令）（電-8・15） \\
中波（AM）放送 & \qtyrange{526.5}{1606.5}{kHz}（局は\qty{9}{kHz}おき） & 総務省 関東総合通信局「ラジオ放送の概要」（電-10・15） \\
FM 放送 & \qtyrange{76.1}{94.9}{MHz}（本文は 76〜95\,MHz） & 総務省 信越総合通信局の報道資料（2020年）（電-15） \\
地上デジタルテレビ放送 & \qtyrange{470}{710}{MHz}（13〜52\,ch） & 総務省の報道資料（2011年）（電-13・15） \\
電波防護指針（一般の環境、\qtyrange{1.5}{300}{GHz}） & 電場 \qty{61.4}{V/m}、磁場の強さ \qty{0.163}{A/m}、電力束密度 \qty{1}{mW/cm^2} & 総務省 九州総合通信局「電波防護のための基準」、電波法施行規則 第21条の4（電-15） \\
東と西の周波数 & 東 \qty{50}{Hz}・西 \qty{60}{Hz}（境はおおむね富士川・糸魚川） & 研究機関の解説・電気学会の資料・会社の記事の3つで照らし合わせた（電-10） \\
\end{longtblr}\par}
{\footnotesize\color{bkGray}電波防護指針の電場÷磁場の強さは $61.4/0.163\approx\qty{377}{\ohm}$（真空の波のインピーダンス $\eta_0$。電-15 ⑦）、電場×磁場の強さは約\qty{10}{W/m^2}＝\qty{1}{mW/cm^2}。\par}

\phbkhead{7　電気の安全と暮らしの目安}
{\footnotesize
\begin{longtblr}{colspec={Q[l,m,wd=40mm] Q[l,m,wd=46mm] X[l,m]}, width=\linewidth, rowhead=1, row{1}={bg=bkVol!14, font=\sffamily\bfseries\footnotesize}, hlines={0.4pt,bkLine}, rowsep=1.4pt, colsep=3.5pt}
名前 & 値 & 出典・注意 \\
体を流れる電流（交流・1秒ほど） & 約\qty{1}{mA}で感じ始め、約\qty{10}{mA}で自分で離れられなくなり、\qty{100}{mA}ほどで心臓が危ない & 病院の資料・電機の会社の解説・厚生労働省「職場のあんぜんサイト」の用語集で照らし合わせた（電-4） \\
漏電遮断器 & \qty{30}{mA}・0.1 秒以内で切るものが多い & 同じ（電-4） \\
電気料金の目安単価 & 31 円/kWh（税込み） & 全国家庭電気製品公正取引協議会の値（2022年に改めた）。照明工業会のガイドで照らし合わせた（電-4） \\
単3アルカリ乾電池の内部抵抗 & 新品で約\qtyrange{0.15}{0.3}{\ohm} & 電池の会社のデータシート（電-5。例題は使った電池で\qty{0.5}{\ohm}とした） \\
テスターの直流電圧の入力抵抗 & 約\qty{10}{M\ohm} & 計測器の会社の技術資料・仕様書（電-5） \\
\end{longtblr}\par}

\phbkhead{8　線の特性インピーダンスと電力の機器}
{\footnotesize
\begin{longtblr}{colspec={Q[l,m,wd=40mm] Q[l,m,wd=46mm] X[l,m]}, width=\linewidth, rowhead=1, row{1}={bg=bkVol!14, font=\sffamily\bfseries\footnotesize}, hlines={0.4pt,bkLine}, rowsep=1.4pt, colsep=3.5pt}
名前 & 値 & 出典・注意 \\
同軸ケーブル（測定器・無線） & \qty{50}{\ohm}（RG-58。内側の導体の半径 \qty{0.45}{mm}、外側の導体の内側の半径 \qty{1.6}{mm}、ポリエチレン。速度係数 0.66） & ケーブルの会社のデータシート2つ（電-3・8・13・14。波-7） \\
同軸ケーブル（テレビ） & \qty{75}{\ohm}（5C-FB、発泡ポリエチレン） & JIS C 3502、ケーブルの会社のカタログ（電-13。波-7） \\
LAN ケーブル（カテゴリー6） & \qty{100}{\ohm} & ケーブルの会社のデータシート（ISO/IEC 11801 を引く）（電-13） \\
USB・HDMI（2本1組の線） & 約\qty{90}{\ohm}・約\qty{100}{\ohm} & 半導体の会社の設計資料（電-13） \\
架空送電線 & 約\qtyrange{300}{400}{\ohm}（1相 約\qty{1}{mH/km}・約\qty{0.011}{\micro F/km}） & 送電線の問題集の数・解説（電-13） \\
半波長アンテナの放射抵抗 & 約\qty{73}{\ohm} & 細い半波長アンテナの式で計算（電-15） \\
送電・配電の電圧 & 送電 500・275・154・66\,kV など、配電 6.6\,kV、家庭 100/200\,V & 百科事典・電機の会社の用語解説・電気学会の資料（2025年の時点。電-11・13） \\
柱上変圧器 & 6600\,V／210-105\,V（単相3線） & JIS C 4304 にそった製品の仕様、入札資料の銘板（電-11・12） \\
配電用の高効率変圧器 & 効率 約99\,\% & 環境省「温室効果ガス削減の対策の手引き」の変圧器の資料（電-12） \\
三相誘導電動機のトップランナー制度 & 2015年4月から IE3（プレミアム効率）、0.75〜\qty{375}{kW} & 雑誌の記事・電機の会社の技報・日本電機工業会の資料（電-12） \\
周波数の目標の幅 & 50/60\,Hz から ±0.2\,Hz 以内（北海道は ±0.3\,Hz） & 電力広域的運営推進機関の報告書・送配電会社のルール・電気学会誌（J-STAGE）を照らし合わせた（電-16） \\
東西の周波数変換設備・北海道と本州の直流連系 & 210万\,kW（2021年）・90万\,kW（2019年、250\,kV） & 送配電会社・電気学会誌・論文誌（J-STAGE）を照らし合わせた（電-16） \\
揚水発電 & 約2700万\,kW（2022年）、総合効率 約70\,\% & 国の資料の抜粋と調査報告を照らし合わせた（電-16。熱-16） \\
日本の発電電力量 & 9911億\,kWh、うち再生可能エネルギー 23.1\,\%（2024年度） & 資源エネルギー庁のエネルギー需給実績の確報（抜粋で照らし合わせた）（電-16） \\
\end{longtblr}\par}

\phbkhead{9　この巻で使う定数}
{\footnotesize
\begin{longtblr}{colspec={Q[l,m,wd=40mm] Q[l,m,wd=46mm] X[l,m]}, width=\linewidth, rowhead=1, row{1}={bg=bkVol!14, font=\sffamily\bfseries\footnotesize}, hlines={0.4pt,bkLine}, rowsep=1.4pt, colsep=3.5pt}
名前 & 値 & 出典・注意 \\
真空の光速 $c$ & \qty{299792458}{m/s}（決めた値） & SI。本文は $3.00\times10^8$\,m/s（電-13・14・15） \\
電気素量 $e$ & \qty{1.602176634e-19}{C}（決めた値） & SI（2019年から）。本文は $1.60\times10^{-19}$\,C（電-1） \\
真空の誘電率 $\varepsilon_0$ & \qty{8.8541878188e-12}{F/m} & CODATA 2022。本文は $8.85\times10^{-12}$\,F/m（電-1） \\
クーロンの定数 $k=1/(4\pi\varepsilon_0)$ & \qty{8.98755e9}{N.m^2/C^2} & 計算。本文は $8.99\times10^9$（電-1。数学 8-5） \\
真空の透磁率 $\mu_0$ & \qty{1.25663706127e-6}{N/A^2} & CODATA 2022。$4\pi\times10^{-7}$ との差は相対 $1.3\times10^{-10}$。本文は $4\pi\times10^{-7}$（電-6） \\
真空の波のインピーダンス $\eta_0=\sqrt{\mu_0/\varepsilon_0}$ & \qty{376.730}{\ohm} & CODATA 2022。本文は約\qty{377}{\ohm}（電-15） \\
電子の質量 $m_{\mathrm e}$ & \qty{9.1093837139e-31}{kg} & CODATA 2022（電-1・6） \\
陽子の質量 $m_{\mathrm p}$ & \qty{1.67262192595e-27}{kg} & CODATA 2022（電-1） \\
ファラデー定数 & \qty{96485.33212}{C/mol} & 決めた値の積（$e$×アボガドロ定数）。本文は文字にしない（電-4。熱-10） \\
太陽定数（地球の所の強さ） & \qty{1361}{W/m^2} & IAU 2015 Resolution B3（電-15。熱-13・波-7） \\
\end{longtblr}\par}
{\footnotesize\color{bkGray}定数は NIST の CODATA 2022 の値（physics.nist.gov）。2019年の SI の改定で、電気素量 $e$ を決めた値にしたので、$\mu_0$ と $\varepsilon_0=1/(\mu_0c^2)$ は、はかって決める値になりました（電-6 ③）。\par}

\phbkhead{10　単位と換算}
{\footnotesize
\begin{longtblr}{colspec={Q[l,m,wd=34mm] Q[l,m,wd=34mm] X[l,m]}, width=\linewidth, rowhead=1, row{1}={bg=bkVol!14, font=\sffamily\bfseries\footnotesize}, hlines={0.4pt,bkLine}, rowsep=1.4pt, colsep=3.5pt}
量 & SI の単位 & ほかの単位との関係（単位つきの1で換算する。2-10） \\
電荷 & C（クーロン）＝A·s & 電子1個は $-1.60\times10^{-19}$\,C \\
電流 & A（アンペア）＝C/s & \qty{1}{mA}＝$10^{-3}$\,A \\
電圧・電位 & V（ボルト）＝J/C & 電子ボルト 1\,eV＝$1.60\times10^{-19}$\,J \\
電場 & V/m＝N/C & ― \\
抵抗・コンダクタンス & Ω（オーム）＝V/A、S（ジーメンス）＝1/Ω & 抵抗率 Ω·m、導電率 S/m \\
静電容量 & F（ファラド）＝C/V & \qty{1}{\micro F}＝$10^{-6}$\,F、1\,pF＝$10^{-12}$\,F \\
インダクタンス & H（ヘンリー）＝V·s/A＝Wb/A & 磁気抵抗の単位 1/H＝A/Wb \\
磁束・磁場 & Wb（ウェーバ）＝V·s、T（テスラ）＝Wb/m$^2$ & 1\,T＝$10^4$\,G（ガウス） \\
磁場の強さ・磁化 & A/m & 1\,Oe（エルステッド）は約\qty{79.6}{A/m} \\
電力 & W（ワット）＝V·A＝J/s & 無効電力 var（バール）、皮相電力 VA（ボルトアンペア） \\
電力量 & J（ジュール） & 1\,kWh＝$3.6\times10^6$\,J（5-12） \\
周波数・角周波数 & Hz（ヘルツ）＝1/s、rad/s & $\omega=2\pi f$。\qty{50}{Hz}は約\qty{314}{rad/s} \\
回転数 & min$^{-1}$（毎分） & 毎分3000回転＝\qty{50}{Hz} \\
レベル & dB（デシベル） & 電力の比：$10\log_{10}$、電圧・電流の比：$20\log_{10}$。2倍で約3\,dB・約6\,dB \\
\end{longtblr}\par}
{\footnotesize\color{bkGray}接頭語：T（テラ、$10^{12}$）、G（ギガ、$10^9$）、M（メガ、$10^6$）、k（キロ、$10^3$）、m（ミリ、$10^{-3}$）、\textmu（マイクロ、$10^{-6}$）、n（ナノ、$10^{-9}$）、p（ピコ、$10^{-12}$）。古い本や外国の本の cgs の単位（ガウス・エルステッドなど）は、ここの換算で SI にしてから使おう。\par}

%% ---------------- なぜ？の道すじ ----------------
%% 第1〜3巻の道すじの終点が電の章だったもの（車はどうして動くの？→電-12、空はなぜ青い？→電-15、燃料電池の電圧→電-4、
%%   スマホの声→電-15）を、この巻の中の道として続けた。章の問い（もくじ）17本のあとに、章の中の身近な問いを足した。
\bkbacktitle{なぜ？の道すじ}
{\small 気になる「なぜ？」から、読む章の順番をたどる路線図です。札の横の数はこの巻のページです。ほかの巻の札（力・熱・波・原）は、物理シリーズの『力学』『熱学』『波動』『原子』の章です。使う数学の章は、下の小さな札です。\par}
\vspace{1mm}
\phwhyroute{送電線は、なぜとても高い電圧で電気を送る？}{電-0}{電-0,電-4,電-11,電-12}{4-9}
\phwhyroute{下じきでこすった髪の毛は、なぜ立ち上がる？}{電-1}{電-1}{4-10,8-5}
\phwhyroute{電気の力は、なぜはなれるほど弱くなる？}{電-2}{電-1,電-2}{8-8,8-11}
\phwhyroute{雷が落ちても、金属の車の中の人が無事なのはなぜ？}{電-3}{電-2,電-3}{8-11}
\phwhyroute{ドライヤーは、電気でどうやって熱をつくる？}{電-4}{熱-1,電-4}{5-12}
\phwhyroute{鳥は、なぜ電線にとまっても感電しない？}{電-5}{電-4,電-5}{3-6}
\phwhyroute{電線に電気を流すと、なぜ方位磁針がふれる？}{電-6}{電-6}{4-14,8-10}
\phwhyroute{磁石にくっついたくぎは、なぜ自分も磁石になる？}{電-7}{電-6,電-7}{}
\phwhyroute{IH調理器は、火がないのになぜなべが熱くなる？}{電-8}{電-4,電-6,電-8}{8-11}
\phwhyroute{充電器をコンセントからぬいても、ランプがしばらく光っているのはなぜ？}{電-9}{電-3,電-4,電-9}{6-4}
\phwhyroute{ラジオは、たくさんの電波からどうやって1つの局を選ぶ？}{電-10}{波-2,電-9,電-10}{6-13}
\phwhyroute{送電線の鉄塔では、なぜ電線が3本ずつ組になっている？}{電-11}{電-10,電-11}{6-9}
\phwhyroute{電気自動車のモーターは、どうやって回る？}{電-12}{電-6,電-8,電-11,電-12}{6-9}
\phwhyroute{テレビのアンテナ線に、なぜふつうの電線を使わない？}{電-13}{波-7,電-10,電-13}{8-13}
\phwhyroute{電気と磁気は、別々のものなの？}{電-14}{電-2,電-6,電-8,電-14}{8-14}
\phwhyroute{電子レンジの窓の網は、何のためにある？}{電-15}{電-8,電-14,電-15}{8-13}
\phwhyroute{停電しないように、電気の量をどうやって合わせている？}{電-16}{電-11,電-12,電-16}{}
\phwhyroute{車はどうして動くの？}{電-12}{力-2,力-6,熱-7,電-12}{4-14,5-5}
\phwhyroute{空は、なぜ青い？　夕焼けは、なぜ赤い？}{波-13}{波-2,波-13,電-15}{6-13}
\phwhyroute{スマホの声は、どうやって電波にのせて送る？}{波-14}{波-14,電-10,電-15}{9-12}
\phwhyroute{燃料電池の電圧は、何で決まる？}{熱-10}{熱-10,電-4}{8-3}
\phwhyroute{雷は、どうやって始まるのか？}{電-2}{電-1,電-2,電-3}{}
\phwhyroute{冬にドアノブでパチッとくるのはなぜ？}{電-1}{電-1,電-3}{}
\phwhyroute{スマホの画面は、どうやって指の場所を知る？}{電-3}{電-3}{}
\phwhyroute{電気料金は、何に払っている？}{電-4}{電-4}{5-12}
\phwhyroute{温度計は、どうやって温度を電気で読む？}{電-4}{熱-1,電-4}{}
\phwhyroute{スマホは、どうやって向きを知る？}{電-6}{電-6}{}
\phwhyroute{地球の磁場は、何がつくっている？}{電-7}{電-6,電-7}{}
\phwhyroute{自転車のライトは、どうやって電気をつくる？}{電-8}{電-8}{5-15}
\phwhyroute{IC カードは、電池がないのになぜ動く？}{電-8}{電-8,電-10}{}
\phwhyroute{カメラのフラッシュは、どうやって強い光を一気に出す？}{電-9}{電-3,電-9}{6-4}
\phwhyroute{変圧器は、どうやって電圧を変える？}{電-12}{電-8,電-12}{}
\phwhyroute{エネルギーは、どこを通って豆電球へ届く？}{電-14}{電-4,電-14}{4-14}
\phwhyroute{電波は、どうやってアンテナから出ていく？}{電-15}{電-10,電-13,電-15}{}
\phwhyroute{光は、なぜ何もない宇宙を進める？}{電-15}{波-5,電-14,電-15}{8-13}

%% ---------------- チェック表 ----------------
\bkbacktitle{できるようになったことチェック表}
{\small 各章の「ふりかえり」を集めた表です。章を学び終えたら1回め、次の巻に進む前に2回めのチェックをしよう。自信がない所は、ページの数字の所へもどろう。\par}\vspace{2mm}
% 自動生成（tools/mkchecklist.py p4）。各章の「ふりかえり」の項目を集めたもの（第13回）
\begingroup\footnotesize
\def\bkcb{\tikz[baseline=-0.5ex]\draw[\bkTF, line width=0.7pt, rounded corners=0.5pt] (0,-0.3em) rectangle (0.8em,0.5em);}
\begin{longtblr}{colspec={Q[c,m,wd=11mm] X[l,m] Q[c,m,wd=10mm] Q[c,m,wd=9mm] Q[c,m,wd=9mm]}, width=\linewidth, rowhead=1, colsep=3pt,
  row{1}={bg=bkVol!14, font=\sffamily\bfseries\scriptsize}, hlines={0.4pt,bkLine}, rowsep=0.7pt}
章 & できるようになったこと & ページ & 1回め & 2回め \\
{\sffamily\bfseries\color{\bkTT}電-0} & 発電所から家のドライヤーまで、電気が通る道すじと、とちゅうで働く機械（発電機・変圧器・送電線）を説明できる。 & \pageref{ph:電-0:sec1} & \bkcb & \bkcb \\
 & 何を無視したかを書いてから、電流・電力・電線で失われる電力を、7つの段で求められる。 & \pageref{ph:電-0:ex3} & \bkcb & \bkcb \\
 & 身近な「なぜ？」の答えが、この巻のどの章にあるかを探せる。 & \pageref{ph:電-0:sec5} & \bkcb & \bkcb \\
{\sffamily\bfseries\color{\bkTT}電-1} & こすると物が電気を帯びるわけを、電子が移ることで説明できる。 & \pageref{ph:電-1:sec1} & \bkcb & \bkcb \\
 & クーロンの法則と重ね合わせで、いくつかの電荷のあいだの力と、電場を求められる。 & \pageref{ph:電-1:sec3} & \bkcb & \bkcb \\
 & 一様な電場の中の電荷の動きを、力学の放物運動と同じ考えで求められる。 & \pageref{ph:電-1:sec8} & \bkcb & \bkcb \\
 & 線・円環・平面に広がった電荷の電場を積分で求め、電気双極子の電場が距離の3乗に反比例して弱まることを示せる。 & \pageref{ph:電-1:sec7} & \bkcb & \bkcb \\
{\sffamily\bfseries\color{\bkTT}電-2} & ガウスの法則で、球・長い直線・広い平面の電荷の電場を求められる。 & \pageref{ph:電-2:sec4} & \bkcb & \bkcb \\
 & 電位を使って、電荷を運ぶ仕事や、加速された電子の速さを求められる。 & \pageref{ph:電-2:sec5} & \bkcb & \bkcb \\
 & 等電位面と電気力線の関係を読み、一様な電場で $V=Ed$ を使える。 & \pageref{ph:電-2:sec7} & \bkcb & \bkcb \\
 & ガウスの法則の微分形と $\vec E=-\mathrm{grad}\,\phi$ を、ベクトル解析で説明できる。 & \pageref{ph:電-2:sec8} & \bkcb & \bkcb \\
 & ラプラス方程式を表計算で解き、電場が集まる所を示せる。 & \pageref{ph:電-2:sec11} & \bkcb & \bkcb \\
{\sffamily\bfseries\color{\bkTT}電-3} & コンデンサーにためた電気で、豆電球や LED を光らせられる。 & \pageref{ph:電-3:sec1} & \bkcb & \bkcb \\
 & 導体の中の電場が0になるわけと、静電遮へいを説明できる。 & \pageref{ph:電-3:sec3} & \bkcb & \bkcb \\
 & 平行板の静電容量・つなぎ方・エネルギーを求められる。 & \pageref{ph:電-3:sec4} & \bkcb & \bkcb \\
 & 分極と電束密度で、同軸・球の静電容量を導ける。 & \pageref{ph:電-3:sec8} & \bkcb & \bkcb \\
 & 境界条件・仮想変位・鏡像法を使って、電場や力を求められる。 & \pageref{ph:電-3:sec9} & \bkcb & \bkcb \\
{\sffamily\bfseries\color{\bkTT}電-4} & オームの法則・直列と並列・電力 $P=VI$ で、回路を調べられる。 & \pageref{ph:電-4:sec3} & \bkcb & \bkcb \\
 & 抵抗率から抵抗を求め、電子の流れとジュール熱を説明できる。 & \pageref{ph:電-4:sec5} & \bkcb & \bkcb \\
 & 抵抗の温度変化、電池の内部抵抗、キルヒホッフの法則を使える。 & \pageref{ph:電-4:sec7} & \bkcb & \bkcb \\
 & $\vec j=\sigma\vec E$・連続の式から、オームの法則と電流則を導ける。 & \pageref{ph:電-4:sec6} & \bkcb & \bkcb \\
{\sffamily\bfseries\color{\bkTT}電-5} & 乾電池と豆電球の直列・並列で、電流と明るさがどう変わるかを説明できる。 & \pageref{ph:電-5:sec2} & \bkcb & \bkcb \\
 & 電流の分かれ方・ブリッジ・電流計と電圧計・コンデンサーをふくむ回路を、キルヒホッフの法則で計算できる。 & \pageref{ph:電-5:sec1} & \bkcb & \bkcb \\
 & 網目電流法・節点電位法で回路を行列の式にし、テブナンの定理や Python で解ける。 & \pageref{ph:電-5:sec3} & \bkcb & \bkcb \\
{\sffamily\bfseries\color{\bkTT}電-6} & 電流の周りに輪の形の磁場ができることを方位磁針で確かめ、その向きを右ねじで言える。 & \pageref{ph:電-6:sec2} & \bkcb & \bkcb \\
 & まっすぐな電流・円い電流・ソレノイドがつくる磁場と、磁場の中の電流・荷電粒子・コイルが受ける力を求められる。 & \pageref{ph:電-6:sec4} & \bkcb & \bkcb \\
 & ビオ・サバールの法則とアンペールの法則（積分形）から磁場を導き、$\mathrm{rot}\,\vec B=\mu_0\vec j$ と $\mathrm{div}\,\vec B=0$ を読める。 & \pageref{ph:電-6:sec6} & \bkcb & \bkcb \\
 & ホール効果・速度選別器・サイクロトロンのしくみを、ローレンツ力で説明できる。 & \pageref{ph:電-6:sec10} & \bkcb & \bkcb \\
{\sffamily\bfseries\color{\bkTT}電-7} & 磁石に近づけた鉄が磁石になるわけを、小さな磁石の向きがそろう図で説明できる。 & \pageref{ph:電-7:sec2} & \bkcb & \bkcb \\
 & 磁化 $M$・磁場の強さ $H$・比透磁率を使って、鉄心の中の磁場を求められる。 & \pageref{ph:電-7:sec3} & \bkcb & \bkcb \\
 & 磁気回路を電気回路と同じ形で解き、すき間のある電磁石の磁束を求められる。 & \pageref{ph:電-7:sec5} & \bkcb & \bkcb \\
 & 鉄と空気の境目で磁力線がどう曲がるかを、境界条件から求められる。 & \pageref{ph:電-7:sec7} & \bkcb & \bkcb \\
 & $BH$ 曲線の面積からヒステリシス損を、減磁曲線から永久磁石の動作点を求められる。 & \pageref{ph:電-7:sec8} & \bkcb & \bkcb \\
{\sffamily\bfseries\color{\bkTT}電-8} & 磁石をコイルに出し入れすると電流が流れ、動かす向きと速さで、電流の向きと大きさが変わることを、実験で確かめられる。 & \pageref{ph:電-8:sec1} & \bkcb & \bkcb \\
 & ファラデーの法則とレンツの法則で、動く棒や発電機に生まれる電圧（誘導起電力）の大きさと向きを求められる。 & \pageref{ph:電-8:sec3} & \bkcb & \bkcb \\
 & 自己インダクタンス・相互インダクタンスを求め、コイルにたまるエネルギーを計算できる。 & \pageref{ph:電-8:sec6} & \bkcb & \bkcb \\
 & ファラデーの法則の微分形を読み、磁場のエネルギーから電磁石が鉄を引く力を求められる。 & \pageref{ph:電-8:sec9} & \bkcb & \bkcb \\
 & 渦電流と表皮効果で、IH調理器や、うすい板を重ねた鉄心のしくみを説明できる。 & \pageref{ph:電-8:sec12} & \bkcb & \bkcb \\
{\sffamily\bfseries\color{\bkTT}電-9} & スイッチを入れた瞬間と、十分に時間がたったあとの回路を、コンデンサーとコイルの性質から説明できる。 & \pageref{ph:電-9:sec1} & \bkcb & \bkcb \\
 & 電圧則と部品のきまりから RC・RL 回路の式を立て、時定数と、エネルギーの行き先を求められる。 & \pageref{ph:電-9:sec3} & \bkcb & \bkcb \\
 & LC 回路の電気振動の周波数とエネルギーの行き来、RLC 回路の臨界抵抗を求められる。 & \pageref{ph:電-9:sec7} & \bkcb & \bkcb \\
 & ラプラス変換と数値計算で RLC 回路の過渡現象を求め、力学と電気の相似の表を、電気の法則から説明できる。 & \pageref{ph:電-9:sec9} & \bkcb & \bkcb \\
{\sffamily\bfseries\color{\bkTT}電-10} & 家の交流の周波数と実効値の意味を説明し、\qty{100}{V}の交流の山の電圧を求められる。 & \pageref{ph:電-10:sec2} & \bkcb & \bkcb \\
 & コイル・コンデンサーのリアクタンスと位相のずれ、RLC 直列回路のインピーダンスと共振周波数を求められる。 & \pageref{ph:電-10:sec4} & \bkcb & \bkcb \\
 & 電圧・電流を複素数（フェーザ）で表して交流回路を解き、Q 値・フィルター・交流ブリッジを計算できる。 & \pageref{ph:電-10:sec5} & \bkcb & \bkcb \\
{\sffamily\bfseries\color{\bkTT}電-11} & 交流の電力が時間とともにゆれること、家の電気と送電線の電気の違いを、図で説明できる。 & \pageref{ph:電-11:sec2} & \bkcb & \bkcb \\
 & 力率を使って、有効電力・無効電力・皮相電力を計算できる。 & \pageref{ph:電-11:sec4} & \bkcb & \bkcb \\
 & コンデンサーで力率をよくすると、電線の電流と損失が減るわけを説明できる。 & \pageref{ph:電-11:sec5} & \bkcb & \bkcb \\
 & 三相交流の Y 結線・$\Delta$ 結線の電圧と電流、三相の電力を計算できる。 & \pageref{ph:電-11:sec7} & \bkcb & \bkcb \\
 & 三相の電力が時間によらず一定になること、回転磁界ができることを、式で示せる。 & \pageref{ph:電-11:sec9} & \bkcb & \bkcb \\
{\sffamily\bfseries\color{\bkTT}電-12} & 磁石とコイルでモーターが回るわけを、電流が磁場から受ける力で説明できる。 & \pageref{ph:電-12:sec1} & \bkcb & \bkcb \\
 & 直流モーターのトルクと回る速さの関係を、逆起電力を使って説明できる。 & \pageref{ph:電-12:sec4} & \bkcb & \bkcb \\
 & 変圧器の電圧と電流の比を、巻数から求められる。 & \pageref{ph:電-12:sec5} & \bkcb & \bkcb \\
 & 変圧器の鉄心の磁束、損失と効率を計算できる。 & \pageref{ph:電-12:sec6} & \bkcb & \bkcb \\
 & 誘導モーターのすべりとトルク、同期モーターとインバーターのしくみを説明できる。 & \pageref{ph:電-12:sec7} & \bkcb & \bkcb \\
{\sffamily\bfseries\color{\bkTT}電-13} & ケーブルの上を電気の波が進み、先やつなぎ目ではね返ることを、速さ・往復の時間・波長で説明できる。 & \pageref{ph:電-13:sec1} & \bkcb & \bkcb \\
 & 特性インピーダンスの違いから、電圧の反射係数と、はね返る電力の割合を求められる。 & \pageref{ph:電-13:sec1} & \bkcb & \bkcb \\
 & 四端子定数の行列で箱をつなぎ、短い配電線の電圧降下と、変圧器の短絡電流を見積もれる。 & \pageref{ph:電-13:sec3} & \bkcb & \bkcb \\
 & 小さな箱のくさりから電信方程式を導き、波の速さ $1/\sqrt{L'C'}$ と特性インピーダンス $\sqrt{L'/C'}$ を求められる。 & \pageref{ph:電-13:sec6} & \bkcb & \bkcb \\
 & 損失のある線・長い送電線・4分の1波長の線を、$\cosh$・$\sinh$ と四端子定数で扱える。 & \pageref{ph:電-13:sec9} & \bkcb & \bkcb \\
{\sffamily\bfseries\color{\bkTT}電-14} & 充電しているコンデンサーの極板の間を流れる変位電流と、その周りの磁場を計算できる。 & \pageref{ph:電-14:sec3} & \bkcb & \bkcb \\
 & マクスウェル方程式を積分形と微分形で読み、電磁場のエネルギーの流れ（ポインティングベクトル）を計算できる。 & \pageref{ph:電-14:sec7} & \bkcb & \bkcb \\
 & 動く人から見ると電場と磁場が入れかわることを説明し、電磁場を数値で解く方法（FDTD）を表計算で動かせる。 & \pageref{ph:電-14:sec9} & \bkcb & \bkcb \\
{\sffamily\bfseries\color{\bkTT}電-15} & 電波・光・X 線などが、波長（周波数）の違う同じ電磁波であることを知り、$c=f\lambda$ で波長を計算できる。 & \pageref{ph:電-15:sec1} & \bkcb & \bkcb \\
 & 電磁波が電場と磁場の直角にゆれる横波で、$E=cB$ の関係をもち、エネルギーと運動量を運ぶことを計算で確かめられる。 & \pageref{ph:電-15:sec4} & \bkcb & \bkcb \\
 & ゆれる電荷が出す電磁波（ラーモアの式）、物の表面での反射率、金属の網が電波を通さないわけを、マクスウェル方程式から説明できる。 & \pageref{ph:電-15:sec8} & \bkcb & \bkcb \\
{\sffamily\bfseries\color{\bkTT}電-16} & 発電所から家までの電気の道すじと、東と西で周波数が違うわけ・東西をつなぐしくみを説明できる。 & \pageref{ph:電-16:sec1} & \bkcb & \bkcb \\
 & 整流・インバーター・直流送電を、スイッチの切りかえと「平均の値」で説明し、整流した電圧の平均を計算できる。 & \pageref{ph:電-16:ex2} & \bkcb & \bkcb \\
 & つくる電気と使う電気がずれると周波数が変わるわけを、回る発電機の運動エネルギーで説明し、周波数の下がる速さを見積もれる。 & \pageref{ph:電-16:sec2} & \bkcb & \bkcb \\
 & 調速機の働きを入れた周波数の変わり方を、プログラムで計算して読める。 & \pageref{ph:電-16:prog} & \bkcb & \bkcb \\
\end{longtblr}
\endgroup


%% ---------------- つぎの巻へ ----------------
\bkbacktitle{つぎの巻へ}
\newcommand\bknextpill[1]{\bk@pill{#1}}
{\small 電磁気学で、電気と磁気の法則を導き、回路・機械・電波を調べる道具がそろいました。物理シリーズの『原子』では、この巻で「結果だけ使う」とした所（金属の中の電子、磁石のもと、ダイオードとトランジスタ）のしくみと、光がつぶでもあることを学びます。\par}\vspace{3mm}
\begin{tcolorbox}[enhanced, colframe=phVol5, colback=white, boxrule=0.9pt, arc=2.5mm, left=3mm, right=3mm, top=2mm, bottom=2mm,
  title={物理シリーズの『原子』：量子から半導体・原子力へ}, fonttitle=\sffamily\bfseries, colbacktitle=phVol5, coltitle=white]
\small
\begin{itemize}[leftmargin=1.2\zw, itemsep=2pt, topsep=1pt, label={\color{phVol5}●}]
\item \bk@pill{原-2}\ 光はつぶでもある：光電効果と光子（電-15 の電磁波が、つぶとしてもふるまう）
\item \bk@pill{原-7}・\bk@pill{原-8}\ 原子と周期表・エネルギーバンド（電-4 の電気の通しやすさと、電-7 の磁区・強磁性のもと）
\item \bk@pill{原-9}・\bk@pill{原-10}\ pn 接合とトランジスタ（電-4 の熱電対と、電-12 のインバーターで理想の部品として使ったスイッチ）
\item \bk@pill{原-11}\ 質量はエネルギー：相対論と原子核（電-14 ⑨の「電場と磁場は、見る人で入れかわる」の続き）
\end{itemize}
\end{tcolorbox}
\phbkhead{この巻の道具は、この先でこう使う}
{\small この巻の章と、その道具を使う物理シリーズのほかの巻（『原子』）の章です（章のはじめの「道しるべ」の「この章の先に」から自動で集めた表）。\par}\vspace{2mm}
% 自動生成（tools/mknext.py p4）。各章の道しるべ「この章の先に」から、後ろの巻の章を集めたもの（第13回）
\begin{longtblr}{colspec={Q[l,m,wd=58mm] X[l,m]}, width=\linewidth, rowhead=1, row{1}={bg=bkVol!14, font=\sffamily\bfseries\footnotesize},
  hlines={0.4pt,bkLine}, rowsep=1.5pt, cells={font=\footnotesize}}
この巻の章 & この道具を使う章 \\
{\bknextpill{電-14}\ 4つの式にまとめる：マクスウェル方程式} & {\bknextpill{原-11}\ 質量はエネルギー：相対論と原子核} \\
{\bknextpill{電-15}\ 光も電波も同じ波：電磁波} & {\bknextpill{原-2}\ 光はつぶでもある：光電効果と光子\\ \bknextpill{原-11}\ 質量はエネルギー：相対論と原子核} \\
{\bknextpill{電-16}\ 電気を届ける・電気を変える} & {\bknextpill{原-9}\ 一方通行の電気：pn接合とLED・太陽電池\\ \bknextpill{原-10}\ スイッチと増幅：トランジスタと電子回路\\ \bknextpill{原-15}\ 原子炉の中の中性子：核分裂と原子炉} \\
\end{longtblr}


%% ---------------- さくいん ----------------
\bkbacktitle{さくいん}
{\small 数字はページです。その語がはじめて出てくるページや、詳しく説明しているページをのせています。}\par\vspace{2mm}
\raggedcolumns\begin{multicols}{2}\small\setlength{\parskip}{0pt}
\directlua{bk.printidx()}
\end{multicols}

%% ---------------- 奥付 ----------------
%% さくいんのあとのメモのページの柱を「メモ」にする（第1〜3巻と同じ）
\clearpage\bksetrunning{}{メモ}{メモ}%
\bkpadfour
\global\bk@decofalse\bkfolio{false}
\null\vfill
\noindent\begin{minipage}{\linewidth}\small
{\sffamily\bfseries\Large\color{bkDeep}\phSeries　第4巻}\par\vspace{1mm}
{\sffamily\bfseries\large 電磁気学：場から回路・機械・電波へ}\par\vspace{3mm}
{\color{bkVol}\hrule height 1pt}\vspace{3mm}
\begin{tabular}{@{}l@{\hspace{6mm}}l@{}}
編者 & \phEditor \\
 & 委員長　小学校長（まとめ役） \\
 & 委員　数学・電気工学・機械工学・教育学・量子力学・建築工学・\\
 & \phantom{委員　}半導体工学の大学教授、幼稚園教諭、塾講師、\\
 & \phantom{委員　}特別支援学校教諭、心理士、科学解説者、小学校教諭 \\
デザイン & ユニバーサルデザイン担当委員 \\
挿絵 & イラスト担当委員 \\
先に読む本 & 「物理を理解するための数学」全9巻、物理シリーズ第1〜3巻 \\
版 & 第1版 Ver0（作成中）　2026年11月 \\
\end{tabular}\par\vspace{3mm}
{\color{bkVol}\hrule height 1pt}\vspace{2mm}
{\footnotesize この本の問題の答えは、巻末の「こたえ」にあります。詳しい解き方は、別冊の問題集と指導書にのせる予定です。物性値の出典は、巻末の「物性値・定数・単位」の表にあります。}
\end{minipage}
\makeatother
